% 图契约：赫兹振荡源、传播空间与谐振接收环的实验结构；几何、方向和公式与正文同一约定。
\documentclass[tikz,border=5pt]{standalone}
\usepackage{fontspec,xeCJK,amsmath,amssymb}
\IfFontExistsTF{Times New Roman}{\setmainfont{Times New Roman}}{\setmainfont{TeX Gyre Termes}}
\IfFontExistsTF{Heiti SC}{\setCJKmainfont[AutoFakeBold=1.5]{Heiti SC}}{\setCJKmainfont[AutoFakeBold=1.5]{FandolHei-Regular}}
\usetikzlibrary{arrows.meta,calc,patterns}
\definecolor{ink}{HTML}{243745}
\definecolor{blue}{HTML}{275DAB}
\definecolor{green}{HTML}{177C63}
\definecolor{red}{HTML}{BD4935}
\definecolor{orange}{HTML}{AA6C16}
\definecolor{purple}{HTML}{715393}
\begin{document}
\begin{tikzpicture}[x=1cm,y=1cm,font=\fontsize{9.5}{12}\selectfont,
v/.style={-{Latex[length=2mm,width=1.3mm]},line width=.85pt},
axis/.style={-{Latex[length=1.7mm,width=1.1mm]},ink!65,line width=.55pt},
aux/.style={ink!45,dashed,line width=.55pt},
lab/.style={fill=white,inner sep=1.8pt},
card/.style={draw=ink!25,fill=blue!3,rounded corners=2pt,text width=6.1cm,minimum height=.85cm,inner sep=6pt,align=center}]
\path[use as bounding box](-8,-3.6)rectangle(8,3.65);
\node[font=\bfseries\fontsize{10.5}{13}\selectfont,text width=15.1cm,align=center]at(0.00000,3.22000){M7-H27　赫兹振荡源、传播空间与谐振接收环的实验结构};
\begin{scope}[shift={(-4,0)}]
\draw[ink,line width=1.4pt](-2.00000,-1.50000)--(-2.00000,-0.18000);
\draw[ink,line width=1.4pt](-2.00000,0.18000)--(-2.00000,1.50000);
\draw[red,fill=red!20](-2,-0.18)circle(0.12);
\draw[red,fill=red!20](-2,0.18)circle(0.12);
\node[]at(-2.00000,2.00000){振荡源};
\draw[green,line width=1pt](2.6,0)arc(0:340:.8);
\node[]at(1.80000,1.50000){接收环};
\node[]at(2.20000,-1.40000){微小火花间隙};
\draw[ink!60](2.20000,-1.10000)--(2.55000,-0.15000);
\draw[blue,line width=.9pt,]plot coordinates{(-1.30000,0.00000) (-1.27667,0.06959) (-1.25333,0.13782) (-1.23000,0.20337) (-1.20667,0.26496) (-1.18333,0.32139) (-1.16000,0.37157) (-1.13667,0.41452) (-1.11333,0.44940) (-1.09000,0.47553) (-1.06667,0.49240) (-1.04333,0.49970) (-1.02000,0.49726) (-0.99667,0.48515) (-0.97333,0.46359) (-0.95000,0.43301) (-0.92667,0.39401) (-0.90333,0.34733) (-0.88000,0.29389) (-0.85667,0.23474) (-0.83333,0.17101) (-0.81000,0.10396) (-0.78667,0.03488) (-0.76333,-0.03488) (-0.74000,-0.10396) (-0.71667,-0.17101) (-0.69333,-0.23474) (-0.67000,-0.29389) (-0.64667,-0.34733) (-0.62333,-0.39401) (-0.60000,-0.43301) (-0.57667,-0.46359) (-0.55333,-0.48515) (-0.53000,-0.49726) (-0.50667,-0.49970) (-0.48333,-0.49240) (-0.46000,-0.47553) (-0.43667,-0.44940) (-0.41333,-0.41452) (-0.39000,-0.37157) (-0.36667,-0.32139) (-0.34333,-0.26496) (-0.32000,-0.20337) (-0.29667,-0.13782) (-0.27333,-0.06959) (-0.25000,-0.00000) (-0.22667,0.06959) (-0.20333,0.13782) (-0.18000,0.20337) (-0.15667,0.26496) (-0.13333,0.32139) (-0.11000,0.37157) (-0.08667,0.41452) (-0.06333,0.44940) (-0.04000,0.47553) (-0.01667,0.49240) (0.00667,0.49970) (0.03000,0.49726) (0.05333,0.48515) (0.07667,0.46359) (0.10000,0.43301) (0.12333,0.39401) (0.14667,0.34733) (0.17000,0.29389) (0.19333,0.23474) (0.21667,0.17101) (0.24000,0.10396) (0.26333,0.03488) (0.28667,-0.03488) (0.31000,-0.10396) (0.33333,-0.17101) (0.35667,-0.23474) (0.38000,-0.29389) (0.40333,-0.34733) (0.42667,-0.39401) (0.45000,-0.43301) (0.47333,-0.46359) (0.49667,-0.48515) (0.52000,-0.49726) (0.54333,-0.49970) (0.56667,-0.49240) (0.59000,-0.47553) (0.61333,-0.44940) (0.63667,-0.41452) (0.66000,-0.37157) (0.68333,-0.32139) (0.70667,-0.26496) (0.73000,-0.20337) (0.75333,-0.13782) (0.77667,-0.06959) (0.80000,-0.00000)};
\draw[v,blue](-0.80000,1.50000)--(0.60000,1.50000);
\node[]at(-0.10000,1.90000){传播};
\end{scope}
\node[card]at(4.00000,1.25000){$\text{振荡源}\longrightarrow\text{空间传播}\longrightarrow\text{接收环}$};
\node[card]at(4.00000,-1.05000){$\text{接收信号具有传播延迟与谐振选择}$};
\end{tikzpicture}\end{document}
